\chapter{Discussion}
\label{chap:discussion}

The prompt-only baseline outperformed both retrieval-based systems on every
criterion and in every question category. This chapter discusses the main
reasons for this result, what it suggests about retrieval in this setting,
and which limitations matter when interpreting the findings.

\section{Why Prompt-Only Performed Best}
\label{sec:interpretation}

One reason is that the base model already seems to possess strong
mathematical knowledge at the Bac level. Topics such as complex numbers,
integrals, recurrence, and convergence are not rare or highly specialized.
A strong multilingual language model can often solve such problems directly,
without needing external support. This is consistent with prior work showing
that large language models internalize substantial procedural and domain
knowledge during pretraining~\cite{brown2020gpt3}. In this setting,
retrieval was therefore not filling a major knowledge gap.
A second reason is that retrieval can hurt when the context is only partly
relevant. The RAG system was designed to rely closely on retrieved material.
When that material matched the wrong subtopic, or only covered part of what
the student needed, the model was pulled toward a weak answer. The
prompt-only system did not face that problem. It could rely on its internal
knowledge while still being guided by the curriculum and style constraints
encoded in the prompt.
Prompt design itself also mattered. The prompt-only system used a carefully
structured prompt that defined the persona, enforced curriculum boundaries,
stabilized answer format, and included explicit instructions for uncertainty
and refusal. In the current setup, this form of guidance turned out to be
more reliable than grounding the model in retrieved context that was not
always appropriate.

\section{Why RAG and Hybrid Scored Lower}
\label{sec:rag-hybrid-discussion}

The main bottleneck was retrieval quality. At first glance, the retrieval
metrics looked strong: all 19 analyzed queries had best L2 distances below
1.0. Under normal conditions, such values would suggest that the system was
finding good matches.
However, the chapter-level analysis showed that low distance did not
reliably mean topical relevance. In the broader 15-query analysis, only 9
queries retrieved material from the expected chapter. A recurring pattern
was that broad ``Bac avec corrections'' compilations received low distances
for many different queries, even when they were not the most relevant source
for the actual topic.
The qualitative examples illustrate this clearly. In the probability example
(B02), the system retrieved material related to a different variant of the
problem despite obtaining the lowest distance in the evaluation. In the
continuity example (C01), retrieval returned material from the right general
area but not from the exact subtopic the student was asking about. These
cases show that retrieval could look strong numerically while still being
misleading for the tutoring task.
The hybrid engine was intended to reduce this problem by adapting its
behavior to retrieval quality. In practice, it almost never behaved
differently from the RAG system. All evaluated queries fell into the
strong-retrieval range defined by the thresholds, so the parts of the hybrid
design meant to handle weaker retrieval were never activated. This largely
explains why the hybrid and RAG scores were so close overall.
Another issue is that the second retrieval pass over course material was
never triggered. The system therefore relied entirely on correction-style
documents and never switched to theorem-based support during evaluation.
This matters especially for coaching-style questions. When a student asks
for the method or the underlying idea, a worked correction to a different
exercise is often less useful than a short theorem-based explanation. This
likely contributed to the weaker performance on Category~C.

\section{What This Suggests About Retrieval in This Setting}
\label{sec:value-of-retrieval}

These results do not show that retrieval is useless for mathematics
tutoring. They show something narrower: in the present configuration,
retrieval did not provide a net benefit over careful prompt design.
An important point is that the mathematical content itself is not the most
domain-specific part of the problem. A strong language model already knows
much of the underlying mathematics. What is more specific to the Tunisian
Baccalaureate is the expected correction style, the curriculum boundaries,
and the way answers are presented in exam preparation.
In principle, this is exactly where retrieval should help most. The system
should be able to retrieve authentic corrections and use them as stylistic
and methodological examples. But in the current setup, retrieval often
introduced context that was noisy, overly broad, or only partially relevant.
Instead of helping the model produce more Bac-faithful answers, it often
interfered with both style and content.
This interpretation is in line with the broader RAG literature.
Gao~et~al.~\cite{gao2024rag_survey} note that retrieval can reduce
performance when the retrieved context is not of sufficient quality. The
present results provide a concrete example of that risk in a small,
curriculum-specific educational corpus.

\section{Limitations}
\label{sec:limitations}

The first and most important limitation is corpus size and composition. The
evaluation was carried out on a corpus of 1{,}422 chunks derived from 624
source files. For a retrieval-based system, this is still relatively small,
especially when many documents are heterogeneous or span multiple topics.
It is therefore possible that the ranking between systems would change if
the corpus were expanded with more exam sessions, more correction variants,
and more chapter-specific material. The present results should be understood
as evidence about this corpus and this stage of development, not as a
general statement that prompt-only systems will always outperform RAG for
math tutoring.
A second limitation is the use of a single evaluator. The grading was
performed by one experienced Tunisian Bac mathematics teacher, which gives
the evaluation strong domain relevance, but it does not allow measurement of
inter-rater agreement. Different teachers might weigh rigor, pedagogy, and
stylistic adherence somewhat differently.

The question set was also relatively small. Only 15 questions were graded on
the main rubric, and the category sizes were uneven. This is enough to show
clear tendencies, but not enough to support strong statistical claims.
Some category-level differences, especially in the smaller categories,
should therefore be interpreted cautiously.
Another limitation is that the evaluation focused on expert judgment rather
than actual student learning outcomes. A teacher-based evaluation is useful
for assessing correctness, style, and pedagogical quality, but it does not
replace a user study. Whether the best-rated system would also produce the
best learning gains remains an open question.
The hybrid thresholds are another limitation. They were chosen empirically
during development, but the evaluation showed that they were too permissive
for the observed distance range. Since almost all queries were treated as
strong retrieval cases, the hybrid design could not really express its
intended behavior. More careful threshold calibration, or a different
retrieval-quality signal altogether, would be worth exploring in future
work.
Finally, the digitized corpus is not error-free. Although the OCR quality
was generally good, occasional transcription errors may still affect
retrieval-based systems by distorting indexed formulas or retrieved context.
This issue does not affect the prompt-only baseline in the same way, which
gives that baseline a structural advantage in the current comparison.

\section{Comparison with Related Work}
\label{sec:comparison-related-work}

Direct comparison with prior systems is difficult because, to the best of
our knowledge, no earlier work has built a retrieval-based tutor
specifically for Tunisian Bac mathematics.
The closest related work comes from two directions. General RAG systems
usually target factual question answering over large knowledge
bases~\cite{gao2024rag_survey}. The present system differs in an important
way: retrieval is not mainly used to supply missing facts, but to provide
worked examples and correction style that the model can imitate. This makes
the retrieval problem more sensitive to pedagogical fit and document
structure than in standard open-domain question answering.

<<<<<<< HEAD
Research on AI in education has shown growing interest in LLM-based
tutoring, but most systems are not built around a specific national
curriculum or exam format~\cite{kasneci2023chatgpt_education}. In that
sense, this thesis adds evidence that prompt engineering and retrieval do
not have a simple hierarchy. When the base model already knows much of the
mathematics, a well-designed prompt can outperform retrieval if the corpus
is too small or too noisy to provide consistently relevant support.
The hybrid setup is also related in spirit to adaptive retrieval approaches
such as CRAG~\cite{yan2024practical_rag} and Self-RAG~\cite{asai2023self_rag}.
However, the present implementation is much simpler and was tested in a much
smaller fixed corpus. The main lesson is not that adaptive retrieval does
not help, but that lightweight routing alone cannot compensate for retrieval
signals that are overly optimistic about context quality.
=======
\subsection{What This Means for Retrieval in Education}
\label{subsec:value-of-retrieval}

These results do not mean that RAG is useless for math tutoring. They mean
that \textit{with this specific corpus size, embedding model, and generative
model}, retrieval did not provide a net benefit over careful prompt
engineering.

The key insight is about what makes Bac math ``domain-specific.'' The
mathematical content itself (derivatives, integrals, complex numbers) is
universal—any strong LLM knows it. What is specific to Tunisia is the
\textit{redaction style}: the formal conventions (``On a\ldots'',
``Donc\ldots'', ``D'après le théorème de\ldots''), the curriculum boundaries,
and the exam format. This is exactly where retrieved corrections should add
the most value. And yet, the Bac Style Adherence criterion showed one of
the widest gaps \textit{against} the retrieval systems (Prompt-Only 3.73
vs.\ RAG 2.93), suggesting that noisy retrieval actively interfered with
style rather than helping it.

This is consistent with the broader RAG literature.
Gao~et~al.~\cite{gao2024rag_survey} note that a RAG system with poor
retrieval can perform worse than the base model alone, because irrelevant
context misleads the generator. Our results are a concrete instance of this
observation in the educational domain.

\subsection{Limitations}
\label{subsec:limitations}

Several limitations constrain how far these conclusions can be generalized.

\subsubsection{Corpus Size and Composition}
\label{subsec:lim-corpus}

This is arguably the most important limitation. The evaluation was conducted
with a corpus of only 1{,}422 chunks derived from 624 source files. This is
small by RAG standards. The core finding—that Prompt-Only beats RAG—may
not hold if the corpus were significantly expanded.

With more data—more years of Bac exams, more correction variants, more
exercise series, more course material—the retrieval system would have a
better chance of finding \textit{exactly} the right correction for each
query. The chapter mismatch problem (40\% chapter mismatch) could
improve if the corpus contained enough examples per chapter to dominate the
multi-topic compilations in the embedding space. The distance--relevance gap
might narrow as the corpus becomes denser in each topic area.

In other words, the RAG system was tested with a limited corpus, and its
performance is likely sensitive to corpus scale. A substantially larger
corpus could change the relative ranking of the systems. Our results reflect
a snapshot at a particular corpus size, not an inherent limitation of
retrieval-augmented approaches for math tutoring.

\subsubsection{Single Evaluator}

The entire evaluation was conducted by a single teacher. While this
evaluator has deep domain expertise, a single rater cannot capture the full
range of pedagogical preferences. Different teachers might weigh
mathematical rigor versus pedagogical clarity differently, or have varying
tolerance for stylistic deviations. Inter-rater reliability could not be
assessed. The results should be interpreted as one expert's assessment, not
a consensus.

\subsubsection{Small Question Set}

Fifteen graded questions across four categories is a modest sample. The
per-category breakdowns (5, 5, 3, and 2 questions) are too small for
meaningful significance testing. Observed differences—particularly the
Hybrid engine's relative strength on Category~D (2 questions)—are
suggestive, not conclusive.

\subsubsection{Model Dependency}

All three systems use Gemini 2.5~Flash. A weaker model that lacks strong
parametric mathematical knowledge might benefit more from retrieval, which
could reverse the ranking. The results are specific to this model and
cannot be assumed to hold for other language models.

\subsubsection{No Student Evaluation}

We measured answer quality as judged by an expert teacher, not actual
learning outcomes. A system that produces technically correct answers might
not be the most effective for learning; conversely, a system with lower
accuracy but better pedagogical scaffolding might help students more. A
proper user study with Tunisian Bac students was outside the scope of this
thesis.

\subsubsection{Retrieval Threshold Calibration}

The Hybrid engine's distance thresholds (Case~A $\leq 1.2$, Case~B
$\leq 1.6$) were set based on preliminary testing. The evaluation showed
that BGE-M3 produces distances below 1.0 for virtually any query against a
corpus of this size, which effectively collapsed the three-case architecture
into a single case. More aggressive thresholds (e.g., Case~A $\leq 0.6$)
might have activated Cases~B and~C, but without better chapter-match
precision, lower thresholds could also misroute queries. Threshold
calibration remains an open question.

\subsubsection{OCR Noise}

The corpus was digitized using Gemini as a multimodal OCR engine. With an
average CER of 5.4\% and \LaTeX{} math accuracy of 97.1\%, the corpus is
not dominated by errors—but it is not error-free either. A single
corrupted exponent or dropped sign in a retrieved correction can mislead
the generative model's entire solution. This affects only the RAG and
Hybrid systems, giving the Prompt-Only baseline an inherent advantage.


\subsection{Comparison with Related Work}
\label{subsec:comparison-related-work}

Direct comparison with existing systems is difficult because, to our
knowledge, no prior work has built a RAG-based tutoring system specifically
for the Tunisian Baccalaureate. The closest related efforts fall into two
categories.

General-purpose RAG systems~\cite{gao2024rag_survey} typically target
factual question answering over large knowledge bases. Our system is
different: the retrieval target is not facts but \textit{stylistic and
methodological exemplars}—we want to find a worked correction whose
redaction style the model can mimic, not retrieve information the model does
not know.

AI-in-education systems~\cite{kasneci2023chatgpt_education} have explored
LLM-based tutoring but typically without retrieval and without targeting a
specific national exam format. Our work contributes the observation that for
a curriculum with strong stylistic conventions, the choice between retrieval
and prompt engineering is not straightforward: a well-engineered prompt can
outperform retrieval when the model already knows the math and the corpus is
too small to provide consistently good matches.

The adaptive retrieval approaches of CRAG~\cite{yan2024practical_rag} and
Self-RAG~\cite{asai2023self_rag} are architecturally related to our Hybrid
engine. Our results suggest that even with quality-aware routing, retrieval
provides limited benefit when the base model is strong and the corpus is
small—complementing the CRAG and Self-RAG findings, which were
demonstrated on larger-scale tasks where parametric knowledge alone is less
likely to suffice.
>>>>>>> 8b150e6 (Fix chapter-match numbers (9/15=60%) and \section heading in discussion)
